\section{Problems for Stage III}
\label{problems:stage-III}

\begin{enumerate}[label=\textbf{III.\arabic*},leftmargin=*]
  \item Substitute the logarithmic WKB ansatz into
  Eq.~\eqref{eq:wkb-schrodinger} and derive the first four Riccati
  coefficients.  Verify the parity relation between the two branches.

  \item Calculate the Borel transform of the factorially divergent series
  $\sum_{n\geq0}n!\hbar^{n+1}$.  Evaluate its two lateral sums and find their
  nonperturbative difference.

  \item Reduce a simple turning point to the Airy equation.  Derive the local
  connection matrix with the book's ordering of dominant and subdominant
  solutions.

  \item Multiply three generic triangular Stokes matrices and one propagation
  matrix.  Determine which orderings preserve the Wronskian and show why path
  order cannot be replaced by a commutative product.

  \item For a barrier with two real turning points, derive the leading
  outgoing resonance condition and the Gamow width.  State which all-order
  object replaces the classical barrier action.

  \item Starting from the hypergeometric connection formula for the inverted
  Rosen--Morse potential, derive the incoming coefficient and its Gamma-
  function zeros.  Identify the resonant and antiresonant branches.

  \item Let $S(E)=N(E)/\Cincoming(E)$.  Derive its Laurent expansion
  at a simple pole and at a double zero.  Anticipate the Jordan resolvent term
  that will appear in Stage VII.

  \item Draw Stokes graphs for a cubic metastable potential on both sheets of
  $\sqrt{Q}$.  Mark a path that imposes outgoing behavior and record every
  branch and Stokes crossing in order.

  \item Demonstrate explicitly that a simultaneous analytic coordinate
  rotation of turning points, cycles, and integration contour leaves a Voros
  period invariant until a singularity crosses the deformation region.

  \item Implement a transfer-matrix calculation for a piecewise constant
  barrier and compare its pole condition with the WKB connection condition.
  Use the determinant as a numerical error monitor.
\end{enumerate}
